% !TEX root = ../../report.tex

\section{Описание сущностей базы данных}

В соответствии с таблицей~\ref{tbl:entities}, описывающей сущности предметной области, выделяются следующие таблицы:

\begin{itemize}
    \item users --- таблица пользователей;
    \item sellers --- таблица продавцов;
    \item products --- таблица товаров;
    \item categories --- таблица категорий;
    \item product\_categories --- таблица связей товаров и категорий;
    \item addresses --- таблица адресов доставки;
    \item orders --- таблица заказов;
    \item order\_items --- таблица позиций заказов;
    \item reviews --- таблица отзывов;
    \item product\_price\_history --- таблица истории изменения цен.
\end{itemize}

В таблицах~\ref{tbl:t_users}--\ref{tbl:t_price_history} приведены столбцы и их ограничения для каждой из перечисленных таблиц.

\begin{table}[H]
    \caption{Описание столбцов таблицы users}
    \label{tbl:t_users}
    \begin{tabular}[]{|c|c|c|c|}
        \hline
        \textbf{Столбец} & \textbf{Тип данных} & \textbf{Ограничения} & \textbf{Описание} \\ \hline
        id & BIGSERIAL & \makecell{NOT NULL \\ PRIMARY KEY} & идентификатор \\ \hline
        email & VARCHAR(255) & \makecell{NOT NULL \\ UNIQUE} & электронная почта \\ \hline
        password\_hash & VARCHAR(255) & NOT NULL & хеш пароля \\ \hline
        full\_name & VARCHAR(255) & NOT NULL & полное имя \\ \hline
        phone & VARCHAR(20) & --- & номер телефона \\ \hline
        role & VARCHAR(20) & \makecell{NOT NULL \\ CHECK} & \makecell{роль: buyer, seller,\\admin, analyst} \\ \hline
        created\_at & TIMESTAMPTZ & \makecell{NOT NULL \\ DEFAULT now()} & дата создания \\ \hline
        deleted\_at & TIMESTAMPTZ & --- & \makecell{дата мягкого\\удаления} \\ \hline
    \end{tabular}
\end{table}

\begin{table}[H]
    \caption{Описание столбцов таблицы sellers}
    \label{tbl:t_sellers}
    \begin{tabular}[]{|c|c|c|c|}
        \hline
        \textbf{Столбец} & \textbf{Тип данных} & \textbf{Ограничения} & \textbf{Описание} \\ \hline
        id & BIGSERIAL & \makecell{NOT NULL \\ PRIMARY KEY} & идентификатор \\ \hline
        user\_id & BIGINT & \makecell{NOT NULL \\ UNIQUE \\ FOREIGN KEY} & \makecell{ссылка на\\пользователя} \\ \hline
        company\_name & VARCHAR(255) & NOT NULL & название компании \\ \hline
        description & TEXT & --- & описание \\ \hline
        rating & NUMERIC(3,2) & --- & \makecell{рейтинг\\(пересчитывается\\триггером)} \\ \hline
        created\_at & TIMESTAMPTZ & \makecell{NOT NULL \\ DEFAULT now()} & дата создания \\ \hline
    \end{tabular}
\end{table}

\begin{table}[H]
    \caption{Описание столбцов таблицы products}
    \label{tbl:t_products}
    \begin{tabular}[]{|c|c|c|c|}
        \hline
        \textbf{Столбец} & \textbf{Тип данных} & \textbf{Ограничения} & \textbf{Описание} \\ \hline
        id & BIGSERIAL & \makecell{NOT NULL \\ PRIMARY KEY} & идентификатор \\ \hline
        seller\_id & BIGINT & \makecell{NOT NULL \\ FOREIGN KEY} & \makecell{ссылка на\\продавца} \\ \hline
        name & VARCHAR(255) & NOT NULL & название товара \\ \hline
        description & TEXT & --- & описание \\ \hline
        price & NUMERIC(12,2) & \makecell{NOT NULL \\ CHECK $> 0$} & цена \\ \hline
        stock\_quantity & INTEGER & \makecell{NOT NULL \\ CHECK $\geq 0$ \\ DEFAULT 0} & \makecell{количество\\на складе} \\ \hline
        reserved\_quantity & INTEGER & \makecell{NOT NULL \\ CHECK $\geq 0$ \\ DEFAULT 0 \\ CHECK $\leq$\\stock\_quantity} & \makecell{количество\\в резерве} \\ \hline
        rating & NUMERIC(3,2) & --- & \makecell{рейтинг\\(пересчитывается\\триггером)} \\ \hline
        created\_at & TIMESTAMPTZ & \makecell{NOT NULL \\ DEFAULT now()} & дата создания \\ \hline
        deleted\_at & TIMESTAMPTZ & --- & \makecell{дата мягкого\\удаления} \\ \hline
    \end{tabular}
\end{table}

\begin{table}[H]
    \caption{Описание столбцов таблицы categories}
    \label{tbl:t_categories}
    \begin{tabular}[]{|c|c|c|c|}
        \hline
        \textbf{Столбец} & \textbf{Тип данных} & \textbf{Ограничения} & \textbf{Описание} \\ \hline
        id & BIGSERIAL & \makecell{NOT NULL \\ PRIMARY KEY} & идентификатор \\ \hline
        parent\_id & BIGINT & FOREIGN KEY & \makecell{ссылка на\\родительскую\\категорию} \\ \hline
        name & VARCHAR(255) & \makecell{NOT NULL \\ UNIQUE} & название категории \\ \hline
        description & TEXT & --- & описание \\ \hline
    \end{tabular}
\end{table}

\begin{table}[H]
    \caption{Описание столбцов таблицы product\_categories}
    \label{tbl:t_product_categories}
    \begin{tabular}[]{|c|c|c|c|}
        \hline
        \textbf{Столбец} & \textbf{Тип данных} & \textbf{Ограничения} & \textbf{Описание} \\ \hline
        product\_id & BIGINT & \makecell{NOT NULL \\ FOREIGN KEY} & \makecell{ссылка на\\товар} \\ \hline
        category\_id & BIGINT & \makecell{NOT NULL \\ FOREIGN KEY} & \makecell{ссылка на\\категорию} \\ \hline
        \multicolumn{4}{|c|}{PRIMARY KEY (product\_id, category\_id)} \\ \hline
    \end{tabular}
\end{table}

\begin{table}[H]
    \caption{Описание столбцов таблицы addresses}
    \label{tbl:t_addresses}
    \begin{tabular}[]{|c|c|c|c|}
        \hline
        \textbf{Столбец} & \textbf{Тип данных} & \textbf{Ограничения} & \textbf{Описание} \\ \hline
        id & BIGSERIAL & \makecell{NOT NULL \\ PRIMARY KEY} & идентификатор \\ \hline
        user\_id & BIGINT & \makecell{NOT NULL \\ FOREIGN KEY} & \makecell{ссылка на\\пользователя} \\ \hline
        city & VARCHAR(100) & NOT NULL & город \\ \hline
        street & VARCHAR(255) & NOT NULL & улица \\ \hline
        house & VARCHAR(20) & NOT NULL & дом \\ \hline
        apartment & VARCHAR(20) & --- & квартира \\ \hline
        zip\_code & VARCHAR(20) & NOT NULL & почтовый индекс \\ \hline
        is\_default & BOOLEAN & \makecell{NOT NULL \\ DEFAULT false} & \makecell{признак основного\\адреса} \\ \hline
        created\_at & TIMESTAMPTZ & \makecell{NOT NULL \\ DEFAULT now()} & дата создания \\ \hline
    \end{tabular}
\end{table}

\begin{table}[H]
    \caption{Описание столбцов таблицы orders}
    \label{tbl:t_orders}
    \begin{tabular}[]{|c|c|c|c|}
        \hline
        \textbf{Столбец} & \textbf{Тип данных} & \textbf{Ограничения} & \textbf{Описание} \\ \hline
        id & BIGSERIAL & \makecell{NOT NULL \\ PRIMARY KEY} & идентификатор \\ \hline
        user\_id & BIGINT & \makecell{NOT NULL \\ FOREIGN KEY} & \makecell{ссылка на\\покупателя} \\ \hline
        seller\_id & BIGINT & \makecell{FOREIGN KEY \\ (nullable)} & \makecell{ссылка на продавца\\(заполняется при\\оформлении)} \\ \hline
        address\_id & BIGINT & \makecell{FOREIGN KEY \\ (nullable)} & \makecell{ссылка на адрес\\(null для\\корзины-черновика)} \\ \hline
        status & VARCHAR(30) & \makecell{NOT NULL \\ CHECK} & \makecell{статус: draft,\\pending, paid,\\shipped, delivered,\\cancelled} \\ \hline
        total\_amount & NUMERIC(12,2) & \makecell{NOT NULL \\ CHECK $\geq 0$} & итоговая сумма \\ \hline
        created\_at & TIMESTAMPTZ & \makecell{NOT NULL \\ DEFAULT now()} & дата создания \\ \hline
        updated\_at & TIMESTAMPTZ & \makecell{NOT NULL \\ DEFAULT now()} & дата обновления \\ \hline
    \end{tabular}
\end{table}

\begin{table}[H]
    \caption{Описание столбцов таблицы order\_items}
    \label{tbl:t_order_items}
    \begin{tabular}[]{|c|c|c|c|}
        \hline
        \textbf{Столбец} & \textbf{Тип данных} & \textbf{Ограничения} & \textbf{Описание} \\ \hline
        id & BIGSERIAL & \makecell{NOT NULL \\ PRIMARY KEY} & идентификатор \\ \hline
        order\_id & BIGINT & \makecell{NOT NULL \\ FOREIGN KEY} & ссылка на заказ \\ \hline
        product\_id & BIGINT & \makecell{NOT NULL \\ FOREIGN KEY} & ссылка на товар \\ \hline
        quantity & INTEGER & \makecell{NOT NULL \\ CHECK $> 0$} & количество \\ \hline
        price\_at\_purchase & NUMERIC(12,2) & \makecell{NOT NULL \\ CHECK $> 0$} & \makecell{цена на момент\\покупки} \\ \hline
        \multicolumn{4}{|c|}{UNIQUE (order\_id, product\_id)} \\ \hline
    \end{tabular}
\end{table}

\begin{table}[H]
    \caption{Описание столбцов таблицы reviews}
    \label{tbl:t_reviews}
    \begin{tabular}[]{|c|c|c|c|}
        \hline
        \textbf{Столбец} & \textbf{Тип данных} & \textbf{Ограничения} & \textbf{Описание} \\ \hline
        id & BIGSERIAL & \makecell{NOT NULL \\ PRIMARY KEY} & идентификатор \\ \hline
        user\_id & BIGINT & \makecell{NOT NULL \\ FOREIGN KEY} & \makecell{ссылка на\\пользователя} \\ \hline
        product\_id & BIGINT & \makecell{NOT NULL \\ FOREIGN KEY} & ссылка на товар \\ \hline
        rating & SMALLINT & \makecell{NOT NULL \\ CHECK 1..5} & оценка от 1 до 5 \\ \hline
        comment & TEXT & --- & текст отзыва \\ \hline
        created\_at & TIMESTAMPTZ & \makecell{NOT NULL \\ DEFAULT now()} & дата создания \\ \hline
        \multicolumn{4}{|c|}{UNIQUE (user\_id, product\_id)} \\ \hline
    \end{tabular}
\end{table}

\begin{table}[H]
    \caption{Описание столбцов таблицы product\_price\_history}
    \label{tbl:t_price_history}
    \begin{tabular}[]{|c|c|c|c|}
        \hline
        \textbf{Столбец} & \textbf{Тип данных} & \textbf{Ограничения} & \textbf{Описание} \\ \hline
        id & BIGSERIAL & \makecell{NOT NULL \\ PRIMARY KEY} & идентификатор \\ \hline
        product\_id & BIGINT & \makecell{NOT NULL \\ FOREIGN KEY} & ссылка на товар \\ \hline
        old\_price & NUMERIC(12,2) & NOT NULL & старая цена \\ \hline
        new\_price & NUMERIC(12,2) & NOT NULL & новая цена \\ \hline
        changed\_at & TIMESTAMPTZ & \makecell{NOT NULL \\ DEFAULT now()} & дата изменения \\ \hline
        changed\_by & TEXT & \makecell{NOT NULL \\ DEFAULT\\current\_user} & \makecell{инициатор\\изменения} \\ \hline
    \end{tabular}
\end{table}

Помимо первичных ключей и ограничений уникальности, для оптимизации запросов создаются дополнительные индексы. Перечень индексов приведён в таблице~\ref{tbl:indexes}.

\begin{table}[H]
    \caption{Индексы базы данных}
    \label{tbl:indexes}
    \begin{tabular}{|l|l|l|}
        \hline
        \textbf{Индекс} & \textbf{Столбцы} & \textbf{Назначение} \\ \hline
        \makecell[l]{idx\_products\_\\seller\_id} & \makecell[l]{products\\(seller\_id)} & \makecell[l]{выборка товаров\\продавца} \\ \hline
        \makecell[l]{idx\_orders\_\\user\_id} & \makecell[l]{orders\\(user\_id)} & \makecell[l]{выборка заказов\\покупателя} \\ \hline
        \makecell[l]{ux\_orders\_one\_\\draft\_per\_user} & \makecell[l]{orders\\(user\_id)\\WHERE status = 'draft'} & \makecell[l]{единственная\\корзина покупателя} \\ \hline
        \makecell[l]{idx\_orders\_\\status\_created\_at} & \makecell[l]{orders\\(status, created\_at)} & \makecell[l]{фильтрация заказов\\по статусу и дате} \\ \hline
        \makecell[l]{idx\_product\_\\categories\_\\category\_id} & \makecell[l]{product\_categories\\(category\_id)} & \makecell[l]{выборка товаров\\по категории} \\ \hline
        \makecell[l]{idx\_reviews\_\\product\_id} & \makecell[l]{reviews\\(product\_id)} & \makecell[l]{выборка отзывов\\на товар} \\ \hline
        \makecell[l]{idx\_products\_\\id\_not\_deleted} & \makecell[l]{products (id)\\WHERE deleted\_at\\IS NULL} & \makecell[l]{частичный индекс:\\только активные\\товары} \\ \hline
        \makecell[l]{idx\_users\_\\id\_not\_deleted} & \makecell[l]{users (id)\\WHERE deleted\_at\\IS NULL} & \makecell[l]{частичный индекс:\\только активные\\пользователи} \\ \hline
        \makecell[l]{idx\_products\_\\price} & \makecell[l]{products\\(price)} & \makecell[l]{сортировка и\\фильтрация по цене} \\ \hline
    \end{tabular}
\end{table}

Частичные индексы (\texttt{WHERE deleted\_at IS NULL}) обеспечивают эффективный поиск только среди активных записей, исключая мягко удалённые строки из индексного дерева.

\clearpage
